% !TEX root = ../ComputationalOTFnT.tex

\chapter{理論的基礎}
\label{c-continuous}

この章では、最適輸送の基礎を説明する。まず確率ベクトル$(\a,\b)$間の最適マッチングとカップリングに関連する概念を導入し、この計算を離散測度$(\al,\be)$間の輸送へと段階的に一般化し、最後に任意の測度に対する一般的な設定を扱う。初読の際には、これらの最後の細部は省略してもよく、読者は確率ベクトル、すなわちヒストグラム間の計算のみに焦点を当てればよい。それが第~\ref{c-algo-basics}章および第~\ref{c-entropic}章で詳述するアルゴリズムを実装するための唯一の前提条件である。より経験豊富な読者は、任意の測度に適用可能な定式化を検討することで問題のより深い理解に到達し、より高度な問題（例えば、点群の位置を移動させる問題や、連続密度から点がサンプリングされる統計的設定など）に適用できるようになるだろう。

%%%%%%%%%%%%%%%%%%%%%%%%%%%%%%%%%%%%%%%%%%%%
\section{ヒストグラムと測度}

確率単体に属する任意の要素$\a \in \simplex_n$に対して、ヒストグラムおよび確率ベクトルという用語を同義に使用する。
\eq{
	\simplex_n \eqdef \enscond{\a \in \RR_+^n}{ \sum_{i=1}^n \a_i = 1 }.
}
本レビューの大部分は、単体上で最適輸送が誘導する幾何学の研究に専ら焦点を当てている。% A more versatile presentation takes advantage of the formalism of discrete measures, which corresponds to handling both the weights contained in a probability vector and the positions associated with the bins of the histogram (the support of the measure). % This formalism can be used implicitly to handle histograms of arbitrary size $n$ and with possibly varying positions. Lastly, the third and most technical layer of exposition deals with arbitrary measures (\ie which need not be discrete, typically continuous with respect to a base measure).


%%%%%%%%%%%%%%
\begin{rem1}{離散測度}
重み$\a$と位置$x_1,\dots,x_n\in\X$を持つ離散測度は次のように表される。
\eql{\label{eq-discr-meas}
	\al = \sum_{i=1}^n \a_i \de_{x_i},
}
ここで$\de_x$は位置$x$におけるディラック測度であり、直感的には位置$x$に無限に集中した単位質量である。このような測度は、さらに$\a\in\simplex_n$であれば確率測度を表し、より一般にベクトル$\a$のすべての要素が非負であれば正測度を表す。質量のない位置が考慮される退化した場合を避けるため、離散測度を考える際には$\a$のすべての要素が正であると仮定する。
\end{rem1}
%%%%%%%%%%%%%%



%%%%%%%%%%%%%%
\begin{rem2}{一般の測度}
最適輸送の便利な特徴は、離散と連続のいずれか、またはその両方の測度を同一の枠組みで扱えることである。そのために、空間$\X$上のラドン測度の集合$\Mm(\X)$を用いる。この集合の厳密な定義には、$\X$が距離（通常$\dist$と記す）を備えていることが必要である。なぜなら、測度には連続関数（$f \in \Cc(\X)$と記す）に対する積分（テスト）を通じてのみアクセスできるからである。

離散測度$\al$に対する$f \in \Cc(\X)$の積分は和として計算される。
\eq{
	\int_\X f(x) \d\al(x) = \sum_{i=1}^n \a_i f(x_i).
}
%
より一般的な測度、例えば$\X=\RR^\dim$（ここで$\dim \in \NN^*$は次元）上の測度は、ルベーグ測度に関する密度$\d\al(x)=\density{\al}(x)\d x$を持つことがある。これはしばしば$\density{\al} = \frac{\d\al}{\d x}$と記され、次を意味する。
\eq{
	\foralls h \in \Cc(\RR^\dim), \quad
	\int_{\RR^\dim} h(x) \d\al(x) =  \int_{\RR^\dim} h(x) \density{\al}(x) \d x.
}
%
任意の測度$\al \in \Mm(\X)$（密度を持つ必要もなく、ディラック測度の和である必要もない）は、任意の連続関数$f \in \Cc(\X)$に対して積分でき$\int_\X f(x) \d\al(x) \in \RR$が得られるという事実によって定義される。$\X$がコンパクトでない場合は、$f$がコンパクト台を持つか、少なくとも無限遠で$0$に収束することも課す必要がある。
%
したがって、測度はある意味で関数より「正則性が低い」が、超関数（滑らかな関数との双対として定義される）よりは正則性が高い。例えば、ディラック測度の導関数は測度ではない。
%
$\Mm_+(\X)$で$\X$上のすべての正測度の集合を表す。確率測度の集合は$\Mm_+^1(\X)$と記され、任意の$\al \in \Mm_+^1(\X)$が正であり、かつ$\al(\X)=\int_\Xx \d\al = 1$であることを意味する。
%
図~\ref{fig-measures}は、本書で考える、ヒストグラム以外のさまざまなクラスの測度の視覚化を提供する。
\end{rem2}
%%%%%%%%%%%%%%


\begin{figure}[h!]
\centering
\begin{tabular}{@{}c@{\hspace{1mm}}c@{\hspace{1mm}}c@{\hspace{1mm}}c@{}}
\includegraphics[width=.24\linewidth]{measures/discrete-1d.pdf}&
\includegraphics[width=.15\linewidth]{measures/discrete-2d}&
\includegraphics[width=.24\linewidth]{measures/density-1d}&
\includegraphics[width=.22\linewidth]{measures/density-2d}\\
離散 $d=1$ & 離散 $d=2$ & 密度 $d=1$ & 密度 $d=2$
\end{tabular}
\caption{\label{fig-measures}
離散分布$\al = \sum_{i=1}^n \a_i \de_{x_i}$（赤は経験的一様分布$\a_i=1/n$に対応し、青は任意の分布に対応）および密度$\d\al(x)=\density{\al}(x)\d x$（紫）の模式図。1次元および2次元の場合を示す。1次元の離散分布はステムプロット（長さは$\a_i$に等しい）で表示され、2次元では点群を用いて表示される（この場合、半径は$\a_i$に等しいか、より視覚的に正確な表現として面積が$\a_i$に等しい）。
}
\end{figure}





%%%%%%%%%%%%%%%%%%%%%%%%%%%%%%%%%%%%%%%%%%%%
\section{割当問題とモンジュ問題}

コスト行列$(\C_{i,j})_{i \in \range{n}, j \in \range{m}}$が与えられ、$n=m$と仮定すると、最適割当問題は$n$個の要素の置換の集合$\Perm(n)$における全単射$\si$を求める問題であり、次を解く。
\eql{\label{eq-optimal-assignment}
	\umin{\si \in \Perm(n)} \frac{1}{n}\sum_{i=1}^n \C_{i,\si(i)}.
}
素朴には、集合$\Perm(n)$のすべての置換を用いて上記のコスト関数を評価できる。しかし、その集合のサイズは$n!$であり、小さな$n$でさえ巨大である。例えば、$n$がわずか70でもその集合は$10^{100}$を超える要素を持つことを考えよ~\citep{Dantzig1983}。したがって、この問題は置換の集合上でコスト関数を最適化する効率的なアルゴリズムが存在する場合にのみ解くことができ、それが\S\ref{s-auction}の主題である。

\begin{rem}[一意性] 最適割当問題は複数の最適解を持つ場合があることに注意せよ。例えば、$n=m=2$であり、行列$\C$が辺の長さ$1$の2次元正方形の4つの頂点間のペアワイズ距離行列であるとする（図~\ref{fig-non-unique-matching}の左のプロットに示す）。この場合、2つの割当のみが存在し、両方とも最適である。
\end{rem}


\begin{figure}[h!]
\centering
\includegraphics[width=.8\linewidth]{non-unique-optimal-matching/non-unique-optimal-matching}
\caption{\label{fig-non-unique-matching}
左：測度$\alpha$の青い点と測度$\beta$の赤い点はペアワイズに等距離である。したがって、マッチング$\sigma=(1,2)$（実線）と$\sigma=(2,1)$（点線）の両方が最適である。右：モンジュ写像は青い測度$\alpha$を赤い測度$\beta$に対応させることができる。重み$\alpha_i$は各位置のディスクの面積に比例して表示されている。ここでの写像は$T(x_1)=T(x_2)=y_2$、$T(x_3)=y_3$であり、$4\leq i\leq 7$に対して$T(x_i)=y_1$である。
}
\end{figure}


%%%%%%%%%
\begin{rem1}{離散測度間のモンジュ問題}\label{rem-monge1}
離散測度
\eql{\label{eq-pair-discr}
	\al = \sum_{i=1}^n \a_i \de_{x_i}
	\qandq
	\be = \sum_{j=1}^m \b_j \de_{y_j},
}
に対して、\citeauthor{Monge1781}の問題~\citeyearpar{Monge1781}は、各点$x_i$を単一の点$y_j$に対応させ、$\al$の質量を$\be$の質量に押し出す写像を求める。すなわち、写像$T:\{x_1,\dots, x_n\}\rightarrow \{y_1,\dots,y_m\}$は次を満たさなければならない。
\eql{\label{eq-monge-constr}
	\foralls j \in \range{m}, \quad
	\b_j = \sum_{ i : \T(x_i) = y_j } \a_i ,
}
これをコンパクトな形で$\T_\sharp \al = \be$と書く。$\b$のすべての要素が正であるため、この写像は必然的に全射である。
%
この写像は、点$(x,y) \in \X \times \Y$に対して定義されるコスト関数$\c(x,y)$によってパラメータ化される輸送コストを最小化しなければならない。
\eql{\label{eq-monge-discr}
	\umin{\T} \enscond{ \sum_{i} \c(x_i,\T(x_i))  }{ \T_\sharp \al = \be } .
}
このような離散点間の写像は、すべての$x$および$y$が異なると仮定すれば、添字$\si : \range{n} \rightarrow \range{m}$を用いて$j=\si(i)$のように符号化でき、質量保存は次のように書かれる。
\eq{
	\sum_{i \in \si^{-1}(j)} \a_i = \b_j,
}
ここで逆$\si^{-1}(j)$は$j$の原像集合として理解する。$n=m$かつすべての重みが一様、すなわち$\a_i=\b_j=1/n$である特殊な場合には、質量保存制約は$\T$が全単射であることを意味し、$\T(x_i)=y_{\si(i)}$となり、モンジュ問題は最適マッチング問題~\eqref{eq-optimal-assignment}と等価になる。ここでコスト行列は
\eq{
	\C_{i,j} \eqdef \c(x_i,y_j).
}
$n\ne m$の場合、最適性は別として、モンジュ写像はある離散測度から別の離散測度へ存在しないこともあることに注意せよ。これは重みベクトルが両立しない場合に起こり、目標測度がソース測度より多くの点を持つ場合、すなわち$n<m$の場合に常にそうなる。例えば、図~\ref{fig-non-unique-matching}の右のプロットは$\al$から$\be$への（最適な）モンジュ写像を示しているが、$\be$から$\al$へのモンジュ写像は存在しない。
\end{rem1}
%%%%%%%%%


%%%%%%%
\begin{rem2}{押し出し作用素}\label{rem-push-f}
連続写像$\T : \X \rightarrow \Y$に対して、対応する押し出し作用素$\T_\sharp : \Mm(\X) \rightarrow \Mm(\Y)$を定義する。
%
離散測度~\eqref{eq-discr-meas}に対して、押し出し操作は単に測度の台にあるすべての点の位置を移動させることからなる。
\eq{
	\T_{\sharp} \al \eqdef \sum_i \a_i \de_{\T(x_i)}.
}
より一般的な測度に対して、例えば密度を持つ測度に対しては、押し出しの概念は確率測度の空間的な変更（または輸送）を記述する上で基本的な役割を果たす。厳密な定義は以下のとおりである。

\begin{defn}[押し出し]\label{defn-pushfwd}
$\T : \X \rightarrow \Y$に対して、ある$\al \in \Mm(\X)$の押し出し測度$\be = \T_\sharp \al \in \Mm(\Y)$は次を満たす。
\eql{\label{eq-push-fwd}
	\foralls h \in \Cc(\Y), \quad \int_\Y h(y) \d \be(y) = \int_\X h(\T(x)) \d\al(x).
	% = \int_X h(T(x)) \density{\al}(x) \d x %%%% y=T(x)   dy=|T'(x)| dx
	% = \int_Y h(y) \density{\al}(T^{-1}(y)) 1/|T'(x)| \d y
	% = \int_Y h(y) \density{\be}(y) \d y
	%  \density{\al}(T^{-1}(y)) 1/|T'(x)| = \density{\be}(y)
}
同値な表現として、任意の可測集合$B \subset \Y$に対して、次が成り立つ。
\eql{\label{eq-equiv-pushfwd}
	\be(B) = \al( \enscond{x \in \X}{\T(x) \in B} ) = \al( T^{-1}(B) ).
}
$\T_\sharp$は正値性と全質量を保存するので、$\al \in \Mm_+^1(\X)$ならば$\T_\sharp \al \in \Mm_+^1(\Y)$となることに注意せよ。
\end{defn}

直感的には、可測写像$T: \X\rightarrow \Y$は単一の点をある可測空間から別の可測空間へ移動させる関数として解釈できる。$T_\sharp$は$T$の拡張であり、$\X$上の確率測度全体を$\Y$上の新しい確率測度へ移動させることができる。作用素$T_\sharp$は、測度$\al$の各要素的質量を写像$T$を適用して$\X$上から$\Y$上の要素的質量へと\emph{押し出す}。押し出し作用素$\T_\sharp : \Mm_+^1(\X) \rightarrow \Mm_+^1(\Y)$は、2つの測度$\al_1,\al_2$に対して$T_\sharp(\al_1+\al_2)=T_\sharp\al_1+ T_\sharp\al_2$という意味で\emph{線形}であることに注意せよ。
\end{rem2}
%%%%%%%

%%%%%%%
\begin{rem2}{多変量密度に対する押し出し}
$\RR^\dim$上の密度$(\density{\al},\density{\be})$を持つ測度に対して（$\T$が滑らかで全単射であると仮定して）、公式~\eqref{eq-push-fwd}における変数変換を明示的に行うと、押し出しが積分公式における変数変換として密度に対して線形に作用することが示される。実際、次が成り立つ。
\eql{\label{eq-pfwd-density}
	\density{\al}(x) = |\det(\T'(x))|  \density{\be}(\T(x)),
}
ここで$\T'(x) \in \RR^{\dim \times \dim}$は$T$のヤコビ行列（$T$の各座標の勾配を取って形成される行列）である。
これは次を意味する。
\eq{
	|\det(\T'(x))| = \frac{ \density{\al}(x) }{ \density{\be}(\T(x)) }.
}
\end{rem2}
%%%%%%%


%%%%%%%
\begin{rem2}{任意の測度間のモンジュ問題}\label{rem-monge2}
モンジュ問題~\eqref{eq-monge-discr}は、2つの任意の確率測度$(\al,\be)$が2つの空間$(\X,\Y)$上に台を持ち、写像$\T : \X \rightarrow \Y$を通じて結ばれる場合に拡張でき、次を最小化する。
\eql{\label{eq-monge-continuous}
	\umin{\T} \enscond{ \int_{\X} \c(x,\T(x)) \d \al(x)  }{  \T_\sharp \al = \be }.
}
制約$\T_\sharp \al = \be$は、$\T$が$\al$の質量を$\be$へ押し出すことを意味し、注意~\ref{rem-push-f}で定義された押し出し作用素を用いている。
\end{rem2}
%%%%%%%


%%%%%%%
\begin{rem2}{押し出しと引き戻し}
測度の押し出し$\T_\sharp$は、関数の引き戻し$\T^\sharp : \Cc(\Y) \rightarrow \Cc(\X)$と混同してはならない。後者は関数間の「ワーピング」に対応し、$g \in \Cc(\Y)$に$\T^\sharp g = g \circ \T$を対応させる線形写像として定義される。押し出しと引き戻しは実際には互いに随伴であり、次の意味においてそうである。
\eq{
	\foralls (\al,g) \in \Mm(\X) \times \Cc(\Y), \quad
	\int_\Y g \d( \T_\sharp\al ) = \int_\X (\T^\sharp g) \d\al.
}
$(\al,\be)$が固定された測度（例えば$\RR^\dim$上のルベーグ測度）に関する密度$(\density{\al},\density{\be})$を持つ場合でも、$\T_\sharp \al$は~\eqref{eq-pfwd-density}におけるヤコビアンの存在により$\T^\sharp \density{\be}$を密度として持つわけではないことに注意せよ。
%
これが、最適輸送を画像レジストレーションに使用する際に注意が必要な理由である。なぜなら、最適輸送は画像ワーピング法としては機能しないからである。
%
図~\ref{fig-push-pull}はこれらの押し出しと引き戻し作用素の区別を示している。
\end{rem2}
%%%%%%%


\begin{figure}[h!]
\centering
\begin{tabular}{@{}c@{\hspace{5mm}}c@{}}
\includegraphics[width=.3\linewidth]{push-pull/push-forward}&
\includegraphics[width=.3\linewidth]{push-pull/pull-back}\\
測度の押し出し & 関数の引き戻し
\end{tabular}
\caption{\label{fig-push-pull}
押し出し作用素$\T_\sharp$（任意の測度を入力として取ることができる）と引き戻し作用素$\T^\sharp$（関数、特に密度に作用する）の比較。
}
\end{figure}



%%%%%%%%%%%%%%
\begin{rem2}{測度と確率変数}\label{rem-meas-random}
	ラドン測度は確率変数の分布を表すものとしても見ることができる。$\X$上の確率変数$X$は実際にはある抽象的な（しばしば不特定の）確率空間$(\Om,\PP)$からの写像$X : \Om \rightarrow \X$であり、その分布$\al$は$\PP(X \in A) = \al(A)=\int_A \d\al(x)$を満たすラドン測度$\al \in \Mm_+^1(\X)$である。
	%
	同値に、これは$\PP$の$X$による押し出し$\al=X_\sharp\PP$である。
	%
	$\T : \X \rightarrow \Y$に対して別の押し出し$\be = \T_\sharp\al$を~\eqref{eq-push-fwd}に従って適用することは、別の確率変数$Y=\T(X) : \om \in \Om \rightarrow \T(X(\om)) \in Y$を定義することと等価であり、$\be$は$Y$の分布となる。
	%
	$Y$からのランダムサンプル$y$の抽出は、$X$から抽出した$x$に対して単に$y=\T(x)$を計算することで達成される。
\end{rem2}
%%%%%%%%%%%%%%


%%%%%%%%%%%%%%%%%%%%%%%%%%%%%%%%%%%%%%%%%%%%
\section{カントロヴィッチ緩和}\label{sec-kantorovich-relaxation}

%%%%%
%\paragraph{Limitations of the Monge Problem}

割当問題、および注意~\ref{rem-monge1}で提示されたモンジュ問題におけるその一般化は、実際の問題に見られるような離散測度を研究するために常に適切とは限らない。実際、割当問題は置換問題として定式化されているため、\emph{同じ}サイズの\emph{一様な}ヒストグラムの比較にのみ使用できる。非一様な重みを持つ離散測度への直接的な一般化はモンジュの押し出し写像の形式を用いて行うことができるが、質量保存制約~\eqref{eq-monge-constr}を満たす実行可能解が存在しない場合にはその定式化も退化し得る（注意~\ref{rem-monge1}の末尾を参照）。さらに、割当問題~\eqref{eq-monge-discr}は組合せ的であり、モンジュ問題~\eqref{eq-monge-continuous}の実行可能集合は、質量保存制約を満たすすべての押し出し測度からなる集合として連続的にパラメータ化されているにもかかわらず、\emph{非凸}である。したがって、いずれも元の定式化のまま取り組むと解くことが困難である。

%%%%%%
% \paragraph{Kantorovich's relaxation} Redundant
\citet{Kantorovich42}の鍵となるアイデアは、輸送の決定論的な性質、すなわちソース点$x_i$が別の点または位置$y_{\sigma_i}$あるいは$T(x_i)$にのみ割り当てられるという事実を緩和することである。カントロヴィッチは代わりに、任意の点$x_i$の質量が複数の位置に分配される可能性を提案する。カントロヴィッチは質量輸送が\emph{決定論的}であるべきという考えから離れ、代わりに\emph{確率論的}な輸送を考え、現在一般に\emph{質量分割}として知られるもの、すなわちソースから複数のターゲットへの分割を可能にする。この柔軟性は、置換$\sigma$や写像$T$の代わりに、カップリング行列$\P  \in \RR_+^{n \times m}$を用いて符号化される。ここで$\P_{i,j}$はビン$i$からビン$j$へ、あるいは離散測度~\eqref{eq-pair-discr}の形式では$x_i$から$y_j$へ流れる質量の量を表す。許容カップリングはモンジュ写像よりもはるかに単純な特徴付けを持つ。
\eql{\label{eq-discr-couplings}
	\CouplingsD(\a,\b) \eqdef \enscond{ \P \in \RR_+^{n \times m} }{
		\P \ones_m = \a \qandq
		\transp{\P} \ones_n = \b
	},
}
ここで次の行列-ベクトル記法を用いた。
\eq{
	\P \ones_m = \left(\sum_j \P_{i,j}\right)_i \in \RR^n
	\qandq
	\transp{\P} \ones_n = \left(\sum_i \P_{i,j}\right)_j \in \RR^m.
}
行列の集合$\CouplingsD(\a,\b)$は有界であり$n+m$個の等式制約で定義されるため、凸多面体（有限個の行列の凸包）である~\citep[\S8.1]{brualdi2006combinatorial}。

%
さらに、モンジュの定式化（図~\ref{fig-non-unique-matching}の右のプロットに示すように）が本質的に非対称であるのに対し、カントロヴィッチの緩和された定式化は常に対称的である。すなわち、カップリング$\P$が$\CouplingsD(\a,\b)$に属するのは$\transp{\P}$が$\CouplingsD(\b,\a)$に属する場合であり、かつその場合に限る。カントロヴィッチの最適輸送問題は次のように読まれる。
\eql{\label{eq-mk-discr}
	\MKD_{\C}(\a,\b) \eqdef
	\umin{\P \in \CouplingsD(\a,\b)}
		\dotp{\C}{\P} \eqdef \sum_{i,j} \C_{i,j} \P_{i,j}.
}
これは線形計画問題であり（第~\ref{c-algo-basics}章を参照）、そのような計画問題の場合に通常そうであるように、その最適解は必ずしも一意ではない。

\begin{rem}[鉱山と工場]\label{rem-kantorovich-primal} カントロヴィッチ問題は、以下の資源配分問題において非常に自然な例示を見出す（\citet{Hitchcock41}も参照）。ある事業者が$n$個の倉庫と$m$個の工場を運営しているとする。各倉庫には工場が適切に稼働するために必要な貴重な原材料が含まれている。より正確には、各倉庫は整数$i$でインデックス付けされ、$\a_i$単位の原材料を含む。これらの原材料はすべて工場に移動させなければならず、工場$j$が適切に機能するためには所定の量$\b_j$が必要である。倉庫$i$から工場$j$へ資源を移送するために、事業者は1単位の資源を位置$i$から位置$j$に移動させるのに$\C_{i,j}$を請求する運送会社を利用できる。運送会社は輸送の独占権を持ち、経済のすべての主体に同一の線形価格体系を適用するものとする。すなわち、$i$から$j$へ$a$単位の資源を輸送するコストは$a\times \C_{i,j}$に等しい。

上記の問題に直面して、事業者は式~\eqref{eq-mk-discr}で記述される線形計画問題を解き、各ペア$i,j$に対して倉庫$i$から工場$j$へ輸送すべき物品の量$\P_{i,j}$を定量化する輸送計画$\P^\star$を得ることを選択する。事業者はその計画を実行するために、運送会社に合計$\dotp{\P^\star}{\C}$を支払う。
\end{rem}

%%%%%%%
\paragraph{カップリングとしての置換行列}

置換$\si\in\Perm(n)$に対して、対応する置換行列を$\P_{\si}$と書く。
	\eql{\label{eq-perm-matrices}
		\foralls (i,j) \in \range{n}^2, \quad
		(\P_{\si})_{i,j} = \choice{
			1/n \qifq j=\si_i, \\
			0 \quad\text{otherwise.}
		}
	}
この場合、次を確認できる。
\eq{\dotp{\C}{\P_{\si}}=\frac{1}{n}\sum_{i=1}^n \C_{i,\si_i},}
これは割当問題~\eqref{eq-optimal-assignment}が、カップリング$\P$を正確に置換行列に制限したカントロヴィッチ問題~\eqref{eq-mk-discr}として書き直せることを示している。
\eq{\umin{\si \in \Perm(n)} \frac{1}{n}\sum_{i=1}^n \C_{i,\si(i)} = \umin{\si \in \Perm(n)} \dotp{\C}{\P_{\si}}.}
次に、置換行列の集合が\citeauthor{birkhoff}多面体$\CouplingsD(\ones_n/n,\ones_n/n)$に真に含まれることが容易に確認できる。実際、任意の置換$\si$に対して$\P_{\si}\ones=\ones_n$かつ$\transp{\P_{\si}}\ones=\ones_n$が成り立つ一方、$\ones_n\transp{\ones_n}/n^2$は有効なカップリングであるが置換行列ではない。したがって、$\dotp{\C}{\P}$の最小値は、置換行列のみを考える場合よりもすべての輸送計画を考える場合のほうが必然的に小さい。
$$ \MKD_{\C}(\ones_n/n,\ones_n/n) \leq \umin{\si \in \Perm(n)} \dotp{\C}{\P_{\si}}.$$

次の命題は、これらの問題が実際には同じ最適値をもたらすこと、すなわち2つの一様測度$\a=\b=\ones_n/n$間のカントロヴィッチ問題~\eqref{eq-mk-discr}を最小化する置換行列が常に見つかることを示す。したがって、カントロヴィッチ緩和は割当問題において\emph{タイト}である。%Note however that some computational algorithm, which are combinatorial in nature, are dedicated to the case of uniform histograms with the same number of points (see Chapter~\ref{c-algo-basics}).
%
図~\ref{fig-matching-kantorovitch}の左側には、この特殊ケースに対応する最適マッチングの2次元の例を示す。
%

\begin{prop}[マッチングに対するカントロヴィッチ]\label{prop-matching-kanto}
	$m=n$かつ$\a=\b=\ones_n/n$ならば、問題~\eqref{eq-mk-discr}の最適解$\P_{\si^\star}$が存在し、それは問題~\eqref{eq-optimal-assignment}の最適置換$\si^\star \in \Perm(n)$に対応する置換行列である。
\end{prop}

\begin{proof}
	\citeauthor{birkhoff}の定理~\citeyearpar{birkhoff}は、$\CouplingsD(\ones_n/n,\ones_n/n)$の端点の集合が置換行列の集合に等しいことを述べている。線形計画の基本定理~\citep[Theorem 2.7]{bertsimas1997introduction}は、空でない多面体上の線形目的関数の最小値は、有限であれば、多面体の端点で達成されることを述べている。
\end{proof}

\begin{figure}[h!]
\centering
\begin{tabular}{@{}c@{\hspace{5mm}}c@{}}
\includegraphics[width=.4\linewidth]{matching-kantorovitch/matching}&
\includegraphics[width=.4\linewidth]{matching-kantorovitch/weighted}
\end{tabular}
\caption{\label{fig-matching-kantorovitch}
%
最適マッチングと一般的なカップリングの比較。$x_i$と$y_j$の間の黒い線分は、表示された最適カップリング$\P_{i,j}$の~\eqref{eq-mk-discr}を解く非ゼロ要素を示す。
%
左：最適マッチング。命題~\ref{prop-matching-kanto}の設定（同じ点数$n=m$の経験測度）に対応。
%
右：これら2つの重み付き点群はマッチングできない。代わりに、カントロヴィッチカップリングを用いて2つの任意の離散測度を対応させることができる。
}
\end{figure}


%%%%%%%
\begin{rem1}{離散測度間のカントロヴィッチ問題}
形式~\eqref{eq-pair-discr}の離散測度$\al,\be$に対して、$\al,\be$の台の点間のすべてのペアワイズコストを行列$\C$に格納する。すなわち$\C_{i,j} \eqdef \c(x_i,y_j)$として、次を定義する。
\eql{\label{eq-kanto-discr}
	\MK_\c(\al,\be) \eqdef \MKD_{\C}(\a,\b).
}
したがって、離散測度間の最適輸送のカントロヴィッチ定式化は、コスト行列$\C$が$\al$と$\be$の台に依存することを除いて、それらに対応する確率重みベクトル$\a,\b$間の問題と同じである。ただし、記法$\MK_\c(\al,\be)$はある状況で有用である。なぜなら、確率重みと台の点の\emph{両方}への依存性を明示するからであり、後者はコスト関数$\c$を通じてのみ考慮される。
\end{rem1}
%%%%%%%


\begin{figure}[h!]
\centering
\begin{tabular}{@{}c@{\hspace{5mm}}c@{\hspace{5mm}}c@{}}
\includegraphics[width=.3\linewidth]{settings/discrete}&
\includegraphics[width=.3\linewidth]{settings/semi-discrete}&
\includegraphics[width=.3\linewidth]{settings/continuous}\\
離散 & 半離散 & 連続
\end{tabular}
\caption{\label{fig-settings}
カントロヴィッチ最適輸送の3つの主要なシナリオで遭遇する入力測度$(\al,\be)$とカップリング$\Couplings(\al,\be)$の模式図。第~\ref{c-algo-semidiscr}章は半離散の設定に特化している。
}
\end{figure}



%%%%%%%
\begin{rem1}{最適割当とカップリングの応用}
最適輸送計画それ自体（カップリング$\P$またはモンジュ写像$\T$が存在する場合のいずれかとして）は、データサイエンス、特に画像処理において多くの応用を見出している。例えば、コントラスト均等化~\citep{delon2004midway}やテクスチャ合成~\citet{gutierrez2017optimal}に使用されてきた。
%
画像科学への最適輸送の応用の重要な部分は、画像マッチング~\citep{zhu2007image,wang2013linear,museyko2009application,li2013novel}、画像融合~\citep{courty2016optimal}、医用画像処理~\citep{wang2011optimal}、形状レジストレーション~\citep{makihara2010earth,lai2014multi,su2015optimal}、および画像透かし~\citep{mathon2014optimal}に関するものである。天体物理学では、最適輸送は初期宇宙の再構成に使用されてきた~\citep{FrischNaturee}。
%
最適輸送は音楽の採譜~\citep{flamary2016optimal}にも使用されており、経済学ではマッチングデータの解釈に多数の応用がある~\citep{galichon2016optimal}。
%
最後に、最適輸送技術を用いて（またはそれに着想を得て）計算された輸送写像は、サンプリング~\citep{reich2013nonparametric,oliver2014minimization}やベイズ推論~\citep{kim2013efficient,el2012bayesian}を行う上でも有用であることを記しておく。
\end{rem1}
%%%%%%%


%%%%%%%
\begin{rem2}{任意の測度間のカントロヴィッチ問題} $\MK_\c$の定義~\eqref{eq-kanto-discr}は、積空間上の同時分布であるカップリング$\pi \in \Mm_+^1(\X \times \Y)$を考えることにより任意の測度に拡張される。離散の場合は、この積測度を$\pi = \sum_{i,j} \P_{i,j} \de_{(x_i,y_j)}$の形に課す特殊な状況である。一般の場合、質量保存制約~\eqref{eq-discr-couplings}は同時確率分布に対する周辺制約として書き直されるべきである。
\eql{\label{eq-coupling-generic}
	\Couplings(\al,\be) \eqdef
	\enscond{
		\pi \in \Mm_+^1(\X \times \Y)
	}{
		P_{\X\sharp} \pi = \al
		\qandq
		P_{\Y\sharp} \pi = \be
	}.
}
ここで$P_{\X\sharp}$と$P_{\Y\sharp}$は射影$P_\X(x,y)=x$と$P_\Y(x,y)=y$の押し出し（定義~\ref{defn-pushfwd}を参照）である。
%
図~\ref{fig-settings}は、異なるクラスの問題（離散測度と密度）に対してこれらのカップリング制約がどのように変換されるかを示している。
%
\eqref{eq-equiv-pushfwd}を用いると、これらの周辺制約は集合$A \subset \X$および$B \subset \Y$に対して$\pi(A \times \Y)=\al(A)$かつ$\pi(\X \times B)=\be(B)$を課すことと等価である。
%
カントロヴィッチ問題~\eqref{eq-mk-discr}は次のように一般化される。
\eql{\label{eq-mk-generic}
	\MK_\c(\al,\be) \eqdef
	\umin{\pi \in \Couplings(\al,\be)}
		\int_{\X \times \Y} \c(x,y) \d\pi(x,y).
}
これは測度の空間上の無限次元線形計画問題である。$(\X,\Y)$がコンパクト空間で$c$が連続であれば、常に解が存在することを容易に示せる。実際、$\Couplings(\al,\be)$は測度の弱位相に対してコンパクトであり（注意~\ref{dfn-weak-conv}を参照）、$\pi \mapsto \int c \d\pi$はこの位相に対する連続関数であり、制約集合は空でない（例えば$\al \otimes \be \in \Couplings(\al,\be)$）。
%
図~\ref{fig-couplings}は~\eqref{eq-mk-generic}を解く離散および連続の最適カップリングの例を示す。
%
図~\ref{fig-couplings-simple}は離散および連続の周辺分布を含む、最適な1次元カップリングの他の例を示す。
%
\todoK{Explain that proof of existence of minimizers (\ie the inf is a min) follows from classical reasoning in variational calculus, especially if the spaces $(\X,\Y)$ are compact (or the support of the input measure) and the cost $\c$ is continuous. Indeed, the topology to use is the weak topology of measure (Definition~\ref{dfn-weak-conv}) and then the function is continuous on a compact. }
\end{rem2}
%%%%%%%

\begin{figure}[h!]
\centering
\begin{tabular}{@{}c@{\hspace{10mm}}c@{}}
\includegraphics[width=.3\linewidth]{couplings/couplings-continuous}&
\includegraphics[width=.3\linewidth]{couplings/couplings-discr}
\end{tabular}
\caption{\label{fig-couplings}
左：密度を持つ2つの1次元測度間の~\eqref{eq-coupling-generic}を解く「連続」カップリング$\pi$。カップリングはモンジュ写像のグラフ$(x,\T(x))$に沿って局在化している（黒で表示）。
%
右：形式~\eqref{eq-pair-discr}の2つの離散測度間の~\eqref{eq-mk-discr}を解く「離散」カップリング$\T$。正のエントリ$\T_{i,j}$は位置$(i,j)$に$\T_{i,j}$に比例する半径の黒いディスクとして表示されている。
}
\end{figure}


\begin{figure}[h!]
\centering
\begin{tabular}{@{}c@{}c@{}c@{}c@{}}
\includegraphics[width=.24\linewidth]{couplings/couplings-1}&
\includegraphics[width=.24\linewidth]{couplings/couplings-2}&
\includegraphics[width=.24\linewidth]{couplings/couplings-3}&
\includegraphics[width=.24\linewidth]{couplings/couplings-4}
\end{tabular}
\caption{\label{fig-couplings-simple}
1次元分布間の最適カップリングの4つの簡単な例。上に写像として（矢印）、下にカップリングとして表現されている。\citet{Levy2017review}に着想を得た。
}
\end{figure}



%%%%%%%%%%%%%%
\begin{rem2}{確率論的解釈}
カントロヴィッチの問題は、注意~\ref{rem-meas-random}に従い、確率変数のプリズムを通じて再解釈できる。実際、問題~\eqref{eq-mk-generic}は次と等価である。
\eql{\label{eq-ot-proba-interpretation}
	\MK_\c(\al,\be) = \umin{(X,Y)} \enscond{ \EE_{(X,Y)}(c(X,Y)) }{ X \sim \al, Y \sim \be },
}
ここで$(X,Y)$は$\X \times \Y$上の確率変数の組であり、$X \sim \al$（あるいは$Y \sim \be$）は$X$（あるいは$Y$）の法則が測度として$\al$（あるいは$\be$）でなければならないことを意味する。組$(X,Y)$の法則は、積空間$\X \times \Y$上の$\pi \in \Couplings(\al,\be)$である。
\end{rem2}
%%%%%%%%%%%%%%



%%%%%%%%%%%%%%%%%%%%%%%%%%%%%%%%%%%%%%%%%%%%
\section{最適輸送の距離的性質}

最適輸送の重要な特徴は、コスト行列が適切な性質を満たす限り、ヒストグラムおよび確率測度間の距離を定義することである。実際、最適輸送は点間の基底距離をヒストグラムまたは測度間の距離に持ち上げる正準的な方法として理解できる。

まず、\citet{RubTomGui00}によって最初に導入された用語を用いて、「基底計量」行列$\C$が固定され、ビン間の置換コストを表し、比較したい複数のヒストグラムにわたって共有される場合を考える。次の命題は、最適輸送がこれらのビン上に台を持つヒストグラム間の有効な距離を提供することを述べている。

\begin{prop}\label{prop-metric-histo}
$n=m$とし、ある$p \geq 1$に対して$\C=\distD^p=(\distD_{i,j}^p)_{i,j} \in \RR^{n \times n}$とする。ここで$\distD \in \RR_+^{n \times n}$は$\range{n}$上の距離であり、すなわち
\begin{enumerate}[label=(\roman*)]
	\item $\distD \in \RR_+^{n \times n}$は対称である。
	\item $\distD_{i,j}=0$であるのは$i=j$の場合に限る。
	\item $\foralls (i,j,k) \in \range{n}^3, \distD_{i,k} \leq \distD_{i,j}+\distD_{j,k}$。
\end{enumerate}
このとき
\eql{\label{eq-wass-p-disc}
	\WassD_p(\a,\b) \eqdef \MKD_{\distD^p}(\a,\b)^{1/p}
}
（$\WassD_p$は$\distD$に依存することに注意）は$\Si_n$上の$p$-ワッサーシュタイン距離を定義する。すなわち$\WassD_p$は対称であり、正であり、$\WassD_p(\a,\b)=0$であるのは$\a = \b$の場合に限り、三角不等式
\eq{
	\foralls \a,\b,\VectMode{c} \in \Si_n, \quad \WassD_p(\a,\VectMode{c}) \leq \WassD_p(\a,\b) + \WassD_p(\b,\VectMode{c}).
}
を満たす。
\end{prop}

\begin{proof}
距離の対称性と正定値性は容易に証明できる。$\C = \distD^p$は対角成分が零であるから、$\WassD_p(\a,\a)=0$であり、対応する最適輸送行列は$\P^\star=\diag(\a)$である。$\distD^p$のすべての非対角要素の正値性により、$\a\ne \b$のとき$\WassD_p(\a,\b)>0$である（この場合、許容カップリングは必然的に対角の外側に非ゼロ要素を持つため）。$\distD^p$の対称性により、$\WassD_p(\a,\b)$自体が対称関数である。


任意の測度に対するワッサーシュタイン距離の三角不等式を証明するために、\citet[Theorem 7.3]{Villani03}は貼り合わせ補題を用い、所定の構造を持つカップリングの存在を強調する。
% Closed forms for such couplings exist in the discrete case which we can readily use.
離散設定では、この貼り合わせカップリングの明示的な構成は簡単である。
%
$\a,\b,\VectMode{c} \in\simplex_n$とする。$\P$と$\Q$をそれぞれ$\a$と$\b$間、および$\b$と$\VectMode{c}$間の輸送問題の2つの最適解とする。
%
$\b$のゼロ座標から生じ得る問題を回避するため、$\tilde\b_j \eqdef \b_j$（$\b_j>0$の場合）、$\tilde{\b}_j\eqdef 1$（それ以外の場合）となるベクトル$\tilde{\b}$を定義し、次を書く。
\eq{
	\SS \eqdef \P \diag(1/\tilde{\b}) \Q \in \RR_+^{n \times n},
}
そして$\SS \in \CouplingsD(\a,\VectMode{c})$であることに注意する。なぜなら
\eq{
	\SS \ones_n  = \P \diag(1/\tilde{\b}) \Q \ones_n = \P (\b / \tilde{\b}) = \P \ones_{\Supp(\b)} = \a,
}
ここで$\ones_{\Supp(\b)}$は$\b_j>0$であるインデックス$j$の位置に1を持ちそれ以外はゼロであるサイズ$n$のベクトルを表し、$\b_j=0$であるような$j$に対しては必然的に$\P_{i,j} = 0$であるから$\P \ones_{\Supp(\b)} = \P \ones = \a$であるという事実を用いている。同様に$\transp{\SS} \ones_n = \VectMode{c}$が確認できる。三角不等式は次から従う。
\begin{align*}
\WassD_p(\a,\VectMode{c})&=\left(\min_{\P\in U(\a,\VectMode{c})}\dotp{\P}{\distD^p}\right)^{1/p} \leq \dotp{\SS}{\distD^p}^{1/p}\\
&= \left(\sum_{ik} \distD^p_{ik}\sum_{j} \frac{\P_{ij}\Q_{jk}}{\tilde{\b}_j}\right)^{1/p} \leq \left(\sum_{ijk} \left(\distD_{ij}+\distD_{jk}\right)^p \frac{\P_{ij}\Q_{jk}}{\tilde{\b}_j}\right)^{1/p} \\
& \leq \left(\sum_{ijk} \distD^p_{ij} \frac{\P_{ij}\Q_{jk}}{\tilde{\b}_j}\right)^{1/p} + \left(\sum_{ijk}\distD^p_{jk} \frac{\P_{ij}\Q_{jk}}{\tilde{\b}_j}\right)^{1/p}.
\end{align*}
第一の不等式は$\SS$の準最適性による。第二は$\distD$の要素に対する三角不等式であり、第三はミンコフスキーの不等式から来る。したがって次を得る。
\begin{align*}
\WassD_p(\a,\VectMode{c})& \leq \left(\sum_{ij} \distD^p_{ij}\P_{ij} \sum_k \frac{\Q_{jk}}{\tilde{\b}_j}\right)^{1/p} + \left(\sum_{jk} \distD^p_{jk} \Q_{jk} \sum_i \frac{\P_{ij}}{\tilde{\b}_j}\right)^{1/p}\\
&= \left(\sum_{ij} \distD^p_{ij}\P_{ij}\right)^{1/p} + \left(\sum_{jk} \distD^p_{jk} \Q_{jk}\right)^{1/p}\\
&= \WassD_p(\a,\b) +\WassD_p(\b,\c),
\end{align*}
これで証明が完了する。
\end{proof}


%%%%%%%
\begin{rem}[$0<p\leq 1$の場合]
$0 < p \leq 1$ならば、$\distD^p$自体が距離であることに注意せよ。これは、$p\geq 1$の場合には$\WassD_p(\a,\b)$が距離であるのに対し、$p \leq 1$の場合には実際には$\WassD_p(\a,\b)^p$が単体上の距離を定義することを意味する。
\end{rem}
%%%%%%%


%%%%%%%
\begin{rem1}{ワッサーシュタイン距離の応用}
最適輸送距離がビン間の基底計量をそのようなビン上のヒストグラム間の距離に自動的に「持ち上げる」という事実は、ヒストグラムを比較するためのコンピュータビジョンや機械学習の応用における有力な手法となっている。
%
これらの分野では、古典的なアプローチとして局所的な特徴量（例えば画像記述子）を「プーリング」し、検索、クラスタリング、または分類を行うために特徴量の経験分布のヒストグラム（いわゆるbag of features）を計算する方法がある。例えば~\citep{oliva2001modeling}を参照。
%
同様の考えに沿って、最適輸送距離は持ち上げられた特徴空間上で信号および画像解析を行うために使用できる~\citep{thorpe2017transportation}。
%
検索とクラスタリングへの応用は画期的な論文~\citep{RubTomGui00}によって開始され、一部の応用（例えばコンピュータビジョン~\citep{pele2008linear,Pele-iccv2009}）に適した閾値行列$\C$に対するより高速なアルゴリズムに伴い新たな応用が生まれた。
%
より最近の応用では、bag-of-wordsに対するearth mover's distanceの使用が強調されており、次元削減~\citep{pmlr-v51-rolet16}やテキスト分類~\citep{kusner2015word,huang2016supervised}を行うため、またはbag-of-wordsを出力する多クラス分類器を訓練するための代替損失を定義するため~\citep{FrognerNIPS}に使用されている。
%
\citet{kolouri2017optimal}は信号処理と機械学習へのそのような応用の最近の概観を提供している。
\end{rem1}
%%%%%%%




%%%%%%%
\begin{rem1}{測度間のワッサーシュタイン距離}
命題~\ref{prop-metric-histo}は、離散である必要のない任意の測度を扱うように一般化できる。

\begin{prop}\label{prop-metric-measure}
$\X=\Y$と仮定し、ある$p \geq 1$に対して$\c(x,y)=\dist(x,y)^p$とする。ここで$\dist$は$\X$上の距離であり、すなわち
\begin{enumerate}[label=(\roman*)]
	\item  $\dist(x,y) = \dist(y,x) \geq 0$、
	\item  $\dist(x,y)=0$であるのは$x=y$の場合に限る、
	\item  $\foralls (x,y,z) \in \X^3, \dist(x,z) \leq \dist(x,y)+\dist(y,z)$。
\end{enumerate}
このとき、$\X$上の$p$-ワッサーシュタイン距離
\eql{\label{eq-defn-wass-dist}
	\Wass_p(\al,\be) \eqdef \MK_{\dist^p}(\al,\be)^{1/p}
}
（$\Wass_p$は$\dist$に依存することに注意）は確かに距離である。すなわち$\Wass_p$は対称であり、非負であり、$\Wass_p(\al,\be)=0$であるのは$\al = \be$の場合に限り、三角不等式
\eq{
	\foralls (\al,\be,\ga) \in  \Mm_+^1(\X)^3, \quad \Wass_p(\al,\ga) \leq \Wass_p(\al,\be) + \Wass_p(\be,\ga).
}
を満たす。
\end{prop}
\begin{proof}
証明は命題~\ref{prop-metric-histo}と同じアプローチに従い、$(\al,\be)$間および$(\be,\ga)$間の最適カップリングを「貼り合わせる」ことで得られる$(\al,\ga)$間のカップリングの存在に依拠する。
\end{proof}
\end{rem1}
%%%%%%%


%%%%%%%
\begin{rem2}{幾何学的直感と弱収束}
ワッサーシュタイン距離$\Wass_p$は多くの重要な性質を持ち、最も重要なのは弱距離であるということ、すなわち台が重ならない特異分布（例えば離散分布）を比較し、2つの分布の台の空間的なずれを定量化できることである。
%
特に、「古典的な」距離（またはダイバージェンス）は離散分布間で定義すらされない（$L^2$ノルムは基底測度に関する密度を持つ連続測度にのみ適用でき、離散$\ell^2$ノルムは位置$(x_i,y_j)$が予め定められた離散集合の値を取る場合にのみ正しく機能する）。これとは鮮明に対照的に、任意の$p> 0$に対して$\Wass_p^p(\de_x,\de_y) = \dist(x,y)$が成り立つ。実際、$\Couplings(\de_x,\de_y)=\{ \delta_{x,y}\}$であることに注意すれば、カントロヴィッチ問題は実行可能解を1つだけ持つので、$\Wass_p^p(\de_x,\de_y)$は必然的に$(\dist(x,y)^p)^{1/p}=\dist(x,y)$である。これは$x \rightarrow y$ならば$\Wass_p(\de_x,\de_y) \rightarrow 0$であることを示す。
%
この性質は、$\Wass_p$が弱収束を定量化する方法であることに対応する。これを以下で定義する。

\begin{defn}[弱収束]\label{dfn-weak-conv}
	コンパクトな定義域$\Xx$上で、$(\al_k)_k$が$\Mm_+^1(\Xx)$において$\al$に弱収束する（$\al_k \rightharpoonup \al$と記す）とは、任意の連続関数$g \in \Cc(\Xx)$に対して$\int_\Xx g \d\al_k \rightarrow \int_\Xx g \d\al$が成り立つことである。
	%
	非コンパクトな定義域では$g$に追加の減衰条件を課す必要がある。
	%
	この弱収束の概念は、確率ベクトルの法則収束に対応する。
\end{defn}

この収束は$\Wass_p(\al_k,\al) \rightarrow 0$と等価であることが示せる~\citep[Theorem 6.8]{Villani09}（非有界距離空間では$p$次までのモーメントの収束も必要）。
\end{rem2}
%%%%%%%


%%%%%%%%%%%%%%
\begin{rem1}{平行移動}
ユークリッド空間$\X=\RR^\dim$上で基底コスト$c(x,y)=\norm{x-y}^2$に対するワッサーシュタイン距離の良い性質は、平行移動を分離できることである。実際、$T_{\tau} : x \mapsto x-\tau$を平行移動作用素とすると、次が成り立つ。
\eq{
	\Wass_2(T_{\tau\sharp} \al,T_{\tau'\sharp} \be)^2  =
	\Wass_2(\al,\be)^2 - 2\dotp{\tau-\tau'}{ \mean_{\al}-\mean_{\be} } + \norm{\tau-\tau'}^2,
}
\todoK{Check this!}
ここで$\mean_{\al} \eqdef \int_\X x \d\al(x) \in \RR^\dim$は$\al$の平均である。
%
特に、これは距離の以下の良い分解を意味する。
\eq{
	\Wass_2(\al,\be)^2  = \Wass_2(\tilde\al,\tilde{\be})^2  + \norm{\mean_{\al}-\mean_{\be}}^2,
}
ここで$(\tilde\al,\tilde{\be})$は「中心化された」零平均測度$\tilde\al=T_{\mean_{\al}\sharp}\al$である。
\end{rem1}
%%%%%%%%%%%%%%
\pagebreak
%%%%%%%
\begin{rem1}{$p = +\infty$の場合}\label{rem-p-inf}
非形式的には、$p \rightarrow +\infty$における$\Wass_p^p$の極限は
\eql{\label{eq-wass-infty}
	\Wass_{\infty}(\al,\be) \eqdef
	\umin{\pi \in \Couplings(\al,\be)}
		\usup{(x,y) \in \Supp(\pi)} d(x,y),
}
であり、ここで上限は$\X^2$上の測度$\pi$に従う本質的上限として理解されるべきである。
%
$p<+\infty$の場合とは対照的に、これは非凸最適化問題であり、数値的に解くことも理論的に研究することも困難である。$\Wass_{\infty}$距離は$(\al,\be)$の台間のハウスドルフ距離と関連している。\S~\ref{sec-hausdorff}を参照されたい。詳細は~\citep{champion2008wasserstein}を参照。
\end{rem1}
%%%%%%%



%%%%%%%%%%%%%%%%%%%%%%%%%%%%%%%%%%%%%%%%%%%%
\section{双対問題}

カントロヴィッチ問題~\eqref{eq-mk-discr}は制約付き凸最小化問題であり、そのため自然にいわゆる双対問題、すなわち制約付き凹最大化問題と対にすることができる。次の基本的な命題は主問題と双対問題の関係を説明する。

\begin{prop}\label{prop-duality-discr}
カントロヴィッチ問題~\eqref{eq-mk-discr}は次の双対問題を認める。
\eql{\label{eq-dual}
	\MKD_\C(\a,\b) =
	\umax{(\fD,\gD) \in \PotentialsD(\C)} \dotp{\fD}{\a} + \dotp{\gD}{\b},
}
ここで許容双対変数の集合は
\eql{\label{eq-feasible-potential}
	\PotentialsD(\C) \eqdef \enscond{
		(\fD,\gD) \in \RR^n \times \RR^m
	}{ \foralls (i,j) \in \range{n} \times \range{m}, \fD \oplus \gD \leq \C }.
}
このような双対変数はしばしば「カントロヴィッチポテンシャル」と呼ばれる。
\end{prop}

\begin{proof}
この結果は、線形計画の強双対性に関するより一般的な結果~\citep[p. 148, Theo. 4.4]{bertsimas1997introduction}の直接的な帰結である。証明の容易な部分、すなわち式~\eqref{eq-dual}の右辺が$\MKD_\C(\a,\b)$の下界であることの確認は、次の節の注意~\ref{rem-duality}で議論する。
%
完全を期して、ラグランジュ双対性を用いて結果を導出する。\eqref{eq-mk-discr}に対応するラグランジアンは次のように読まれる。
\eql{\label{eq-mk-lagr}
	\umin{\P \geq 0} \umax{ (\fD,\gD) \in \RR^n \times \RR^m }
		\dotp{\C}{\P} + \dotp{\a - \P\ones_m}{\fD} + \dotp{\b - \transp{\P} \ones_n}{\gD}.
}
上記の最小と最大を交換する。これは線形計画問題（有限次元の場合）を考えるときに常に可能であり\todoK{donner une ref}、次を得る。
\eq{
	\umax{ (\fD,\gD) \in \RR^n \times \RR^m }
	\dotp{\a}{\fD} + \dotp{\b}{\gD}
	+ \umin{\P \geq 0}
		\dotp{\C - \fD\transp{\ones_m} - \ones_n \transp{\gD}}{\P}.
}
次に注目すべきは
\eq{
	\umin{\P \geq 0} \dotp{\Q}{\P} =
	\choice{
		0 \qifq \Q \geq 0,\\
		-\infty \quad \text{otherwise}
	}
}
であり、したがって制約は$\C - \fD\transp{\ones_m} - \ones_n \transp{\gD} = \C-\fD\oplus \gD \geq 0$と読まれる。
\end{proof}

ラグランジアン~\eqref{eq-mk-lagr}に対する主双対最適性関係により、最適輸送計画の台を特定できる（\S\ref{s-complementary}も参照）。
\eql{\label{eq-mk-pd-rel}
\{(i,j) \in \range{n} \times \range{m} : \P_{i,j}>0\} \subset \enscond{(i,j) \in \range{n} \times \range{m}}{ \fD_i+\gD_j=\C_{i,j} }.
}


\begin{rem}\label{rem-kantorovich-dual}
注意~\ref{rem-kantorovich-primal}でカントロヴィッチ問題に与えた解釈に続き、双対問題の直感的な説明を行う。その設定では、事業者が全体の資源を倉庫から工場へ可能な限り低コストで移動させたいことを思い出されたい。事業者は~\eqref{eq-mk-discr}を解き、$\P^\star$に示された指示に従い、運送会社に$\dotp{\P^\star}{\C}$を支払うことでこれを行える。

\emph{物流のアウトソーシング。} 事業者が線形計画問題~\eqref{eq-mk-discr}を解く計算手段を持っていないと仮定する。事業者はその代わりにそのタスクをベンダーにアウトソーシングすることを決定する。ベンダーは以下の構造の価格体系を選択する。ベンダーは物流タスクを物品の集荷と配送に分割し、各倉庫$i$で1単位の資源を集荷するために集荷価格$\fD_i$を（その単位がどこに送られるかに関係なく）適用し、工場$j$に1単位の資源を配送するために価格$\gD_j$を（その単位がどの倉庫から来るかに関係なく）適用する。合計では、倉庫$i$に正確に$\a_i$単位があり工場$j$に$\b_j$が必要であるから、ベンダーはこの価格体系の結果として事業者の物流問題を解決するために$\dotp{\fD}{\a}+\dotp{\gD}{\b}$の価格を要求する。

\emph{価格の設定。} ベンダーが使用する価格体系は、かなり自然に任意の負の価格を許容することに注意せよ。実際、ベンダーが倉庫に価格ベクトル$\fD$を、工場に価格ベクトル$\gD$を適用する場合、$\fD$のすべてのエントリを任意の数だけ同時に減少させ$\gD$のすべてのエントリを同じ数だけ増加させても合計請求額は変わらない。なぜなら、すべての倉庫の資源の総量は工場に配送されるべきものと等しいからである。つまり、ベンダーは物品の一部を集荷するために事業者に支払うことで極端に良い取引をしているように見せかけることができるが、配送料をそれだけ多く請求することでその損失を補填できる。これを知った上で、ベンダーはそのサービスに対してできるだけ多く請求したいので、ベクトル$\fD$と$\gD$をできるだけ高く設定する。

\emph{価格の確認。} 競合するベンダーがいない場合、事業者はベンダーの価格が妥当であるかを確認する簡単な方法を考えなければならない。一つの方法は、事業者が問題~\eqref{eq-mk-discr}を解いて最も効率的な計画の価格$\MKD_{\C}(\a,\b)$を計算し、ベンダーの提示がその金額以下であるか確認することである。しかし、事業者はそもそもそのような長い計算を行う余裕がないことを思い出されたい。幸い、事業者がベンダーの提案が競争力あるかどうかを確認するはるかに効率的な方法がある。$\fD_i$はベンダーが$i$で1単位を集荷するために請求する価格であり、$\gD_j$は$j$に1単位を配送するための価格であることを思い出されたい。したがって、ベンダーの価格体系は$i$から$j$へ1単位の資源を移送するコストが正確に$\fD_i+\gD_j$であることを意味する。しかし、事業者は$i$から$j$へ1単位を輸送するために運送会社が請求する価格が$\C_{i,j}$であることも知っている。したがって、任意のペア$i,j$に対して合計価格$\fD_i+\gD_j$が$\C_{i,j}$より厳密に大きければ、ベンダーはそのタスクに対して運送会社が請求する公正な価格よりも多く請求しており、事業者はベンダーの提案を拒否すべきである。


\begin{figure}[h!]
\centering
\includegraphics[width=\linewidth]{dual-solutions-overlay/dual_sols_final}
\caption{\label{dual-solutions-overlay}
左のプロットでは、それぞれ青い点と赤い四角で表される2つの離散測度$\al$と$\be$の間の最適輸送問題を考える。これらのマーカーの面積は各位置の重みに比例している。このプロットは二次ユークリッドコストを用いた最適輸送$\P^\star$も表示している。この配置に対応する双対（カントロヴィッチ）ポテンシャル$\fD^\star$と$\gD^\star$も右のプロットに表示されている。$\alpha$の各点に「価格」$\fD_i^\star$があるので（$\gD$と$\beta$についても同様）、その点の色は右のカラーマップを用いて得られた値を表す。これらのポテンシャルは相対価格として解釈でき、最適な輸送計画の下で$\al$の各位置で質量を移動させる個別コスト、あるいは逆に$\be$の任意の点に質量を送るコストを示す。したがって、最適輸送コストは左のすべての弧の長さの二乗にそれらの太さを乗じた和に等しく、あるいは双対定式化を用いて、右のプロットの各マーカーの面積にその値（色で符号化）を乗じた和に等しい。}
\end{figure}


\emph{双対問題としての最適価格。} したがって、事業者はすべてのペア$i,j$に対してベンダーが提供する価格が$\fD_i+\gD_j \leq \C_{i,j}$を満たすことを確認することが利益にかなう。事業者がベンダーの提供した価格ベクトルがこれら$n\times m$個の不等式に適合することを確認したとする。事業者はベンダーの提案が魅力的であると結論できるだろうか？ 手計算で素早く計算すると、事業者は確かにその提案を受け入れることが利益にかなうと結論する。実際、事業者のいかなる輸送計画$\P$もコスト$\dotp{\P}{\C}=\sum_{i,j} \P_{i,j}\C_{i,j}$を持つので、事業者はこれら$n\times m$個の不等式を適用して、任意の輸送計画$\P$（最適なもの$\P^\star$を含む）に対して、周辺制約から次が成り立つと結論できる。
\eq{\begin{aligned}
	\sum_{i,j} \P_{i,j}\C_{i,j} &\geq \sum_{i,j} \P_{i,j} \left(\fD_{i} + \gD_j\right) = \left(\sum_{i} \fD_{i} \sum_{j} \P_{i,j}\right) + \left(\sum_{j} \gD_{j} \sum_{i} \P_{i,j}\right)\\ 	& =\dotp{\fD}{\a}+\dotp{\gD}{\b},
\end{aligned}}
したがって、自分でその仕事を行う\emph{いかなる}試みも必然的にベンダーの価格より高くつくことを観察する。

これを知った上で、ベンダーは$\dotp{\fD}{\a}+\dotp{\gD}{\b}$を最大化する価格の組$\fD,\gD$を見つけなければならないが、提案が受け入れられるためにはすべての$i,j$に対して少なくとも$\fD_i+\gD_j\leq \C_{i,j}$という基本的な不等式を満たさなければならず、これが問題~\eqref{eq-dual}となる。後に\S\ref{s-kantorovich}で示すように、ベンダーが得る最良の価格は、実際に事業者が$\MKD_{\C}(\a,\b)$を計算して得られる最良の可能なコストと正確に等しいことを示すことができる。

図~\ref{dual-solutions-overlay}は同一の輸送問題から生じる主問題と双対問題の解を示している。左では、青い点が倉庫を、赤い点が工場を表し、これらの点の面積は確率重み$\a,\b$を表し、それらの間のリンクは最適輸送を表し、その幅は輸送量に比例する。ベンダーが問題~\eqref{eq-dual}を最適化した結果として得られる最適価格が右に示されている。価格は平均がゼロになるように選ばれている。最も高い相対価格は、図の左下にある孤立した倉庫での物品集荷と、右上に位置する工場への物品配送から生じている。
\end{rem}



%%%%%%%
\begin{rem2}{任意の測度間の双対問題}
	この主双対構成を任意の測度に拡張するためには、測度は自然に連続関数と双対性の関係にあること（測度には連続関数に対する積分を通じてのみアクセスできる）を認識することが重要である。双対性は以下の命題で形式化され、離散測度を扱う場合には命題~\ref{prop-duality-discr}に帰着する。

	\begin{prop}
	次が成り立つ。
	\eql{\label{eq-dual-generic}
		\MK_\c(\al,\be) =
		\usup{(\f,\g) \in \Potentials(\c)}
			\int_\X \f(x) \d\al(x) + \int_\Y \g(y) \d\be(y),
	}
	ここで許容双対ポテンシャルの集合は
	\eql{\label{eq-dfn-pot-dual}
		\Potentials(\c) \eqdef \enscond{
			(\f,\g) \in \Cc(\X) \times \Cc(\Y)
		}{
			\forall (x,y), \f(x)+\g(y) \leq \c(x,y)
		}.
	}
	ここで$(\f,\g)$は連続関数の組であり、離散の場合と同様に「カントロヴィッチポテンシャル」とも呼ばれる。
	\end{prop}

	\todoK{\begin{proof}
	TODO.
	\end{proof}
	}

	離散の場合~\eqref{eq-dual}は、双対ベクトルが連続ポテンシャルのサンプルであることに対応する。すなわち$(\fD_i,\gD_j)=(\f(x_i),\g(y_j))$である。
	%
	主双対最適性条件により最適計画の台を追跡でき、\eqref{eq-mk-pd-rel}は次のように一般化される。
	\eql{\label{eq-mk-pd-rel-cont}
		\Supp(\pi) \subset \enscond{(x,y) \in \X \times \Y}{ \f(x)+\g(y)=\c(x,y) }.
	}

	主問題~\eqref{eq-mk-generic}とは対照的に、\eqref{eq-dual-generic}の解の存在を示すことは自明ではない。なぜなら制約集合$\Potentials(\c)$はコンパクトではなく、最小化すべき関数は強制的でないからである。
	%
	\S\ref{s-c-transform}で詳述する$c$変換の機構を用いると、$c(x,y)=\dist(x,y)^p$（$p\geq 1$の場合）において、最適な$(\f,\g)$は必然的にリプシッツ正則であることを示すことができ、制約をコンパクトなものに置き換えることが可能になる。
\end{rem2}
%%%%%%%

%%%%%%%
\begin{rem2}{制約なし双対問題}\label{rem-uncons-dual}
$\int_\X \d\mu=\int_\Y\d\nu = 1$の場合、制約付き双対問題~\eqref{eq-dual-generic}を制約なしの問題に置き換えることができる。
\eql{\label{eq-dual-uncons}
	\MK_\c(\al,\be) =
		\usup{(\f,\g) \in \Cc(\X) \times \Cc(\Y)}
			\int_\X \f\d\al + \int_\Y \g \d\be
			+ \umin{\X \otimes \Y} ( c - f \oplus g ),
}
ここで$(f\oplus g)(x,y)=f(x)+g(y)$と記した。ここでの最小値は測度$\mu \otimes \nu$に付随する本質的上限として考えるべきであり、すなわち$f$または$g$が$\mu$および$\nu$の零測度の集合上で変更されても変化しない。
%
この代替的な双対定式化はFrancis Bachから指摘を受けたものである。
%
これは主問題~\eqref{eq-mk-generic}に冗長な制約$\int \d\pi=1$を加えることで得られる。
\end{rem2}
%%%%%%%


%%%%%%%
\begin{rem2}{モンジュ--カントロヴィッチの等価性---ブレニエの定理}\label{rem-exist-mongemap}
以下の定理はしばしば\citet{Brenier91}に帰せられ、$\RR^\dim$において$p=2$の場合、2つの入力測度のうち少なくとも一方が密度を持ち、2次のモーメントを持つ測度に対して、カントロヴィッチ問題とモンジュ問題が等価であることを保証する。興味ある読者は、\citet{cuesta1989,ruschendorf1990characterization}によるほぼ同時期に発表された同じ結果の変種、特に~\citep{MR923203}の原論文や\citet{knott1984optimal}による先駆的な結果も参照されたい。

\begin{thm}[ブレニエ]
	$\X=\Y=\RR^\dim$かつ$c(x,y)=\norm{x-y}^2$の場合、2つの入力測度のうち少なくとも一方（$\al$と記す）がルベーグ測度に関する密度$\density{\al}$を持つならば、カントロヴィッチの定式化~\eqref{eq-mk-generic}における最適な$\pi$は一意であり、「モンジュ写像」$\T : \RR^\dim \rightarrow \RR^\dim$のグラフ$(x,\T(x))$上に台を持つ。これは$\pi = (\Id,\T)_{\sharp} \mu$を意味し、すなわち
\eql{\label{eq-brenier-map}
	\foralls h \in \Cc(\X \times \Y), \quad
		\int_{\X \times \Y} h(x,y) \d \pi(x,y) = \int_{\X} h(x,\T(x)) \d\mu(x).
}
さらに、この写像$\T$は凸関数$\phi$の勾配$\T(x) = \nabla\phi(x)$として一意に定義され、$\phi$は$(\nabla\phi)_\sharp \mu=\nu$を満たす（加法定数の違いを除いて）一意の凸関数である。この凸関数は~\eqref{eq-dual-generic}を解く双対ポテンシャル$\f$と$\phi(x)=\frac{\norm{x}^2}{2}-\f(x)$の関係にある。
\end{thm}

\begin{proof}
	証明の主要な要素を概説する。詳細は例えば~\citep{SantambrogioBook}に見出せる。
	%
	$\int c \d\pi = C_{\al,\be} - 2 \int \dotp{x}{y}\d\pi(x,y)$であることに注意する。ここで定数は$C_{\al,\be} = \int \norm{x}^2\d\al(x) + \int \norm{y}^2\d\be(y)$である。\eqref{eq-mk-generic}を解く代わりに、次の問題を考えることができる。
	\eq{
		\umax{\pi \in \Couplings(\al,\be)}
			\int_{\X \times \Y} \dotp{x}{y} \d\pi(x,y),
	}
	その双対問題は
	\eql{\label{eq-brenier-proof-1}
		\umin{(\phi,\psi)}
		\enscond{
			\int_\X \phi \d\al + \int_\Y \psi \d\be
		}{
			\forall (x,y), \quad \phi(x)+\psi(y) \geq \dotp{x}{y}
		}.
	}
	これらの変数と~\eqref{eq-dfn-pot-dual}のものとの関係は
	$(\phi,\psi) = (\frac{\norm{\cdot}^2}{2}-\f,\frac{\norm{\cdot}^2}{2}-\g)$
	である。
	%
	制約を次で置き換えることができる。
	\eql{\label{eq-brenier-proof-leg}
		\foralls y, \quad \psi(y) \geq \phi^*(y) \eqdef \usup{x} \dotp{x}{y} - \phi(x).
	}
	ここで$\phi^*$は$\phi$のルジャンドル変換であり、線形形式の上限として凸関数である（\eqref{eq-legendre}も参照）。\eqref{eq-brenier-proof-2}に現れる目的関数は線形であり積分測度は正であるから、$\psi$に関して明示的に最小化し$\psi=\phi^*$と設定して、制約なし問題を考えることができる。
	\eql{\label{eq-brenier-proof-2}
		\umin{\phi} \int_\X \phi \d\al + \int_\Y \phi^* \d\be;
	}
	\S\ref{sec-c-transforms}および\S\ref{s-c-transform}も参照されたい。そこではこのアイデアがそれぞれ離散設定と一般のコスト$\c(x,y)$に対して適用されている。
	%
	この議論を2回繰り返すことで、$\phi$を$\phi^{**}$（凸関数）で置き換えることができ、したがって~\eqref{eq-brenier-proof-2}において$\phi$が凸であることを課せる。
	%
	条件~\eqref{eq-mk-pd-rel-cont}は、最適な$\pi$が$\enscond{(x,y)}{\phi(x)+\phi^*(y)=\dotp{x}{y}}$上に台を持つことを示し、そのような$y$はルジャンドル変換の最小化~\eqref{eq-brenier-proof-leg}において最適であり、その最適性条件は$y \in \partial \phi(x)$と読まれる。
	%
	$\phi$は凸であるからほとんどいたるところ微分可能であり、$\al$が密度を持つので$\al$-ほとんどいたるところ微分可能でもある。
	%
	これは各$x$に対して、対応する$y$が$\al$-ほとんどいたるところ$y = \nabla\phi(x)$として一意に定義されることを示し、必然的に$\pi = (\Id,\nabla\phi)_\sharp \al$であることを示す。
\end{proof}

この結果は、非特異密度を持つ$\Wass_2$の設定において、モンジュ問題~\eqref{eq-monge-continuous}とそのカントロヴィッチ緩和~\eqref{eq-mk-generic}が等しいこと（緩和がタイトであること）を示している。これは割当の場合~\eqref{prop-matching-kanto}に対する命題~\ref{prop-matching-kanto}の連続版であり、周辺分布が等しく一様のとき最適輸送問題の最小値が置換行列（離散写像）で達成されることを述べている。
%
最適輸送写像が凸関数の勾配でなければならないというブレニエの定理は、1次元以上における増加関数の概念の有用な一般化を提供する。これが、最適輸送が任意の次元での分位関数を定義するために使用でき、それが分位回帰問題への応用に有用である主な理由である~\citep{carlier2016vector}。

この定理は多くの方向に拡張できることにも注意せよ。
%
$\al$が密度を持つという条件は、$\al$がハウスドルフ次元$\dim-1$未満の「小さな集合」（例えば超曲面）に質量を与えないという条件に弱められる。
%
$\c(x,y)=h(x-y)$の形のコスト（$h$は狭義凸関数）も考えることができる。
\end{rem2}
%%%%%%%


%%%%%%%
\begin{rem2}{モンジュ--アンペール方程式}\label{rem:MA}
密度を持つ測度に対して、\eqref{eq-pfwd-density}を用いると、$\phi$は以下のモンジュ--アンペール型方程式を解く（加法定数の違いを除いて）一意の凸関数であることが得られる。
\eql{\label{eq-monge-ampere}
	\det(\partial^2\phi(x))  \density{\be}(\nabla\phi(x)) = \density{\al}(x)
}
ここで$\partial^2\phi(x) \in \RR^{\dim \times \dim}$は$\phi$のヘッセ行列である。モンジュ--アンペール作用素$\det(\partial^2\phi(x))$は非線形退化ラプラシアンとして理解できる。小さな変位の極限において$\phi=\Id + \epsilon\psi$とすると、滑らかな写像に対して次のようにラプラシアン$\Delta$が線形化として回復される。
\eq{
	\det(\partial^2\phi(x)) = 1 + \epsilon \Delta \psi(x) + o(\epsilon).
}
%
凸性制約は$\det(\partial^2\phi(x)) \geq 0$を強制し、この方程式が解を持つために必要である。
%
モンジュ--アンペール方程式の理論的解析、特にその解の正則性に関する膨大な文献がある。例えば~\citep{gutierrez2016monge}を参照されたい。興味ある読者は\citet{caffarelli2003monge}のレビュー論文を参照されたい。
%
主要な困難は、一般的には解が滑らかである必要はなく、入力測度が任意の場合（例えばディラック質量）にはアレクサンドロフ解の理論に頼る必要があることである。
%
固定された右辺$f$に対するモンジュ--アンペール方程式$\det(\partial^2 \phi(x)) = f(x)$のより単純な場合には多くの解法が提案されてきた。例えば~\citep{benamou2016monotone}およびそこでの参考文献を参照されたい。特に、異方的凸関数の捕捉には特別な注意が必要であり、通常の有限差分は不正確になり得る。
%
最適輸送において$f$が実際には$\nabla \phi$に依存する場合、方程式~\eqref{eq-monge-ampere}の離散化と境界条件は、\citep{benamou2014numerical}およびそこでの参考文献で概説される技術的課題をもたらす。
%
不動点反復に基づく関連する解法が画像レジストレーションに適用されていることにも注意されたい~\citep{haker2004optimal}。
\end{rem2}
%%%%%%%


%%%%%%%%%%%%%%%%%%%%%%%%%%%%%%%%%%%%%%%%%%%%
\section{特殊ケース}\label{sec:specialcases}

一般に、最適輸送距離の計算は数値的に困難である。\S\S \ref{c-algo-basics}、\ref{c-entropic}、および\ref{c-dynamic}でさまざまな数値解法を詳述する前に、まず最適輸送問題の解決が比較的容易な特殊な有利なケースを概観する。


%%%%%%%
\begin{rem}[2値コスト行列と1-ノルム]\label{rem-binary}
コスト行列$\C$が対角上で0、それ以外で$1$である場合、すなわち$\C=\ones_{n\times n}-\Identity_n$のとき、$\a$と$\b$間の1-ワッサーシュタイン距離はそれらの差の1-ノルムに等しいこと、$\MKD_\C(\a,\b)=\norm{\a-\b}_1$が容易に確認できる。
\end{rem}
%%%%%%%

%%%%%%%
\begin{rem1}{クロネッカーコスト関数と全変動}
上記の注意~\ref{rem-binary}に加えて、$c(x,y)$が$x=y$のとき0、$x\ne y$のとき1である場合に、この結果が任意の測度に拡張されることも容易に確認できる。2つの離散測度$\al$と$\be$間の最適輸送距離はそれらの全変動距離に等しい（例~\ref{exmp-tv}も参照）。
\end{rem1}
%%%%%%%



%%%%%%%
\begin{rem1}{1次元の場合---経験測度}\label{rem-1d-empir}
ここでは$\X=\RR$とする。$\al = \frac{1}{n}\sum_{i=1}^n \de_{x_i}$および$\be = \frac{1}{n}\sum_{j=1}^n \de_{y_j}$と仮定し、（一般性を失うことなく）点が順序付けられていると仮定する。すなわち$x_1 \leq x_2 \leq \cdots \leq x_n$かつ$y_1 \leq y_2 \leq \cdots \leq y_n$とすると、次の簡単な公式が成り立つ。
\eql{\label{eq-1d-empirical}
	\Wass_p(\al,\be)^p = \frac{1}{n} \sum_{i=1}^n |x_i-y_i|^p,
}
すなわち局所的には（点がすべて異なると仮定すれば）、$\Wass_p(\al,\be)$は$\al$と$\be$の順序付けられた値の2つのベクトル間の$\ell^p$ノルムである。この主張は局所的にのみ有効であり、値の一部が変化するたびに順序（およびそれらのベクトル表現）が変わり得るという意味でそうである。この公式は注意~\ref{rem-1d-ot-generic}で詳述されるより一般的な設定の簡単な帰結である。
%
図~\ref{fig-1d-discrete}の上段は、同じ点数の経験測度間の1次元輸送写像を示している。下段は、この単調写像が任意の離散測度にどのように一般化されるかを示している。


この1次元計算を活用して、\citet{delon-circle}が示したように円上の最適輸送も効率的に計算することが可能である。
%
コストが距離の凹関数である場合、特に$p<1$の場合、最適輸送計画の挙動は大きく異なるが、効率的な解法も存在することに注意せよ~\citep{delon-concave}。
\end{rem1}
%%%%%%%


\begin{figure}[h!]
\centering
\includegraphics[width=.65\linewidth]{1d-discrete/1d-schematic}
\caption{\label{fig-1d-discrete}
1次元最適カップリング：各矢印$x_i \rightarrow y_j$は最適カップリングにおける非ゼロの$\P_{i,j}$を示す。
%
上段：同じ点数の経験測度（最適マッチング）。
下段：一般的な場合。
%
これは単調再配列に対応する。$x_i \leq x_{i'}$が$\P_{i,j} \neq 0, \P_{i',j'} \neq 0$を満たすならば、必然的に$y_j \leq y_{j'}$である。
}
\end{figure}


% \bar x_{\si_1(i)} --> \bar y_{\si_2(i)}  k --> \si_2(\si_1^{-1}(k))
%%%%%%%
\begin{rem1}{ヒストグラム均等化}
1次元の最適輸送はヒストグラム均等化に使用でき、グレースケール画像のパレットの正規化への応用がある。図~\ref{fig-hist-eq}を参照されたい。この場合、$(\bar x_i)_i$と$(\bar y_j)_j$は2つの入力画像のすべてのピクセルの灰色レベル（$0$が黒、$1$が白、およびそれらの間のすべての値）を予め定められた順序（すなわち列方向）で列挙したものである。各画像のピクセル数が同じで$n\times m$に等しいと仮定すると、これらの色レベルをソートすることで注意~\ref{rem-1d-empir}のように$x_i = \bar x_{\si_1(i)}$および$y_j = \bar y_{\si_2(j)}$が定義される。ここで$\si_1, \si_2 : \{1,\ldots,nm\} \rightarrow \{1,\ldots,nm\}$は置換であり、$\si \eqdef \si_2 \circ \si_1^{-1}$が2つの離散分布間の最適割当となる。画像処理の応用では、$(\bar y_{\si(i)})_i$は$\bar x$の均等化されたバージョンの色値を定義し、その経験分布は$\bar y$のものと正確に一致する。その画像の均等化されたバージョンは、その$nm$次元ベクトルをサイズ$n\times m$の画像として折り返すことで復元できる。また、$t \in [0,1] \mapsto (1-t) \bar x_{i} + t \bar y_{\si(i)}$は元の画像と均等化されたものとの間の補間を定義し、そのピクセルの経験分布は入力のものの間の変位補間（\eqref{eq-displacement-1d-cumul}で定義）である。
\end{rem1}
%%%%%%%

\begin{figure}[h!]
\centering
\begin{tabular}{@{}c@{\hspace{1mm}}c@{\hspace{1mm}}c@{\hspace{1mm}}c@{\hspace{1mm}}c@{\hspace{1mm}}@{}}
\includegraphics[width=.19\linewidth]{hist-eq/equalized-1}&
\includegraphics[width=.19\linewidth]{hist-eq/equalized-3}&
\includegraphics[width=.19\linewidth]{hist-eq/equalized-5}&
\includegraphics[width=.19\linewidth]{hist-eq/equalized-7}&
\includegraphics[width=.19\linewidth]{hist-eq/equalized-9}\\
\includegraphics[width=.19\linewidth]{hist-eq/equalize-1}&
\includegraphics[width=.19\linewidth]{hist-eq/equalize-3}&
\includegraphics[width=.19\linewidth]{hist-eq/equalize-5}&
\includegraphics[width=.19\linewidth]{hist-eq/equalize-7}&
\includegraphics[width=.19\linewidth]{hist-eq/equalize-9}\\
$t=0$ & $t=0.25$& $t=0.5$& $t=.75$& $t=1$
\end{tabular}
\caption{\label{fig-hist-eq}
画像処理のためのヒストグラム均等化。$t$はヒストグラム間の変位補間をパラメータ化する。
}
\end{figure}


%%%%%%%
\begin{rem2}{1次元の場合---一般的な場合}\label{rem-1d-ot-generic}
$\RR$上の測度$\al$に対して、$\RR$から$\rightarrow [0,1]$への累積分布関数を次のように導入する。
\eql{\label{eq-cumul-defn}
	\foralls x \in \RR, \quad \cumul{\al}(x) \eqdef \int_{-\infty}^x \d\al,
}
およびその擬逆関数$\cumul{\al}^{-1} : [0,1] \rightarrow \RR \cup \{-\infty\}$
\eql{\label{eq-pseudo-inv-cum}
	\foralls r \in [0,1], \quad \cumul{\al}^{-1}(r) = \umin{x} \enscond{x \in \RR \cup \{-\infty\} }{ \cumul{\al}(x) \geq r }.
}
%
この関数は$\alpha$の一般化分位関数とも呼ばれる。任意の$p \geq 1$に対して、次が成り立つ。
\eql{\label{eq-wass-cumul}
	\Wass_p(\al,\be)^p = \norm{ \cumul{\al}^{-1} - \cumul{\be}^{-1} }_{L^p([0,1])}^p = \int_0^1 | \cumul{\al}^{-1}(r) - \cumul{\be}^{-1}(r) |^p \d r.
}
これは、写像$\al \mapsto \cumul{\al}^{-1}$を通じて、ワッサーシュタイン距離が$L^p$ノルムを備えた線形空間と等距離であること、あるいは同値に、実数直線上の測度に対するワッサーシュタイン距離がヒルベルト的距離であることを意味する。
これにより1次元最適輸送の幾何学は非常に単純になるが、高次元のそれとは大きく異なる。高次元ではヒルベルト的ではなく、これは命題~\ref{prop-negative-definite}および\S\ref{sec-non-embeddability}でより一般的に議論される。
%
$p=1$の場合、さらに単純な公式が成り立つ。
\begin{align}\label{eq-w1-1d}
	\Wass_1(\al,\be) &= \norm{ \cumul{\al} - \cumul{\be} }_{L^1(\RR)} =
	\int_\RR | \cumul{\al}(x) - \cumul{\be}(x) | \d x \\
	&= \int_\RR \abs{ \int_{-\infty}^x \d(\al-\be) } \d x,
\end{align}
これは$\Wass_1$がノルムであることを示す（任意の次元への一般化については\S\ref{sec-w1-eucl}を参照）。
%
$\T_\sharp \al=\be$を満たす最適モンジュ写像$\T$は次で定義される。
\eql{\label{eq-OT-map-1d}
 	\T = \cumul{\be}^{-1} \circ \cumul{\al}.
}
図~\ref{fig-1d-ot}は累積関数を用いた1次元最適輸送の計算を示している。変位補間も表示されており、\eqref{eq-displacement-1d-cumul}で詳述されるように計算される。注意~\ref{rem-bary-1d}も参照されたい。1次元における最適輸送の性質の詳細な調査については、\cite[Chapter 2]{SantambrogioBook}を参照されたい。
\end{rem2}
%%%%%%%


\newcommand{\MyFigCumulMeas}[1]{\includegraphics[width=.33\linewidth]{1d-cumulative/#1}}
\newcommand{\MyFigCumulCum}[1]{\includegraphics[width=.24\linewidth]{1d-cumulative/#1}}
\begin{figure}[ht!]
\centering
\begin{tabular}{@{}c@{\hspace{1mm}}c@{\hspace{1mm}}c@{}}
\MyFigCumulMeas{input-mu}&
\MyFigCumulMeas{input-nu}&
\MyFigCumulMeas{interp-bary}\\
$\mu$ & $\nu$ & ${ (t\T+(1-t)\Id)_\sharp \mu}$
\end{tabular}
%%%%
\begin{tabular}{@{}c@{\hspace{2mm}}c@{\hspace{2mm}}c@{\hspace{2mm}}c@{}}
\MyFigCumulCum{cumul}&
\MyFigCumulCum{icumul}&
\MyFigCumulCum{transports}&
\MyFigCumulCum{interp-cumul}\\
$(\cumul{\al},\cumul{\be})$ &
$(\cumul{\al}^{-1},\cumul{\be}^{-1})$ &
$(T,T^{-1})$ &
$(1-t)\cumul{\al}^{-1}+t\cumul{\be}^{-1}$
\end{tabular}
\caption{\label{fig-1d-ot}
\eqref{eq-OT-map-1d}で詳述される累積関数を用いた、2つの1次元測度間の最適輸送と変位補間の計算。
}
\end{figure}

%%%%%%%
\begin{rem2}{ガウス分布間の距離}\label{rem-dist-gaussians}
$\al = \Nn(\mean_\al,\cov_\al)$と$\be = \Nn(\mean_\be,\cov_\be)$が$\RR^\dim$上の2つのガウス分布であるとき、次の写像
\eql{\label{eq-transport-Bures}T:x\mapsto \mean_\be + A(x-\mean_\al),}
ここで
$$A=\cov_\al^{-\tfrac{1}{2}}\Big(\cov_\al^{\tfrac{1}{2}}\cov_\be\cov_\al^{\tfrac{1}{2}}\Big)^{\tfrac{1}{2}}\cov_\al^{-\tfrac{1}{2}}=\transp{A},$$
が$T_\sharp \rho_\al = \rho_\be$を満たすことを示すことができる。実際、変数変換の公式~\eqref{eq-pfwd-density}が満たされることに注意すればよい。なぜなら
$$
\begin{aligned}\rho_\be(T(x))&=\det(2\pi\cov_\be)^{-\tfrac{1}{2}} \exp(-\dotp{T(x)-\mean_\be}{\cov_\be^{-1}(T(x)-\mean_\be)})\\
&= \det(2\pi\cov_\be)^{-\tfrac{1}{2}} \exp(-\dotp{ x-\mean_\al}{\transp{A}\cov_\be^{-1}A(x-\mean_\al)}) \\
&= \det(2\pi\cov_\be)^{-\tfrac{1}{2}} \exp(-\dotp{ x-\mean_\al}{\cov_\al^{-1}(x-\mean_\al)}),
\end{aligned}$$
であり、$T$は線形写像であるから
$$|\det T'(x)|= \det A = \left(\frac{\det\cov_\be}{\det\cov_\al}\right)^{\tfrac{1}{2}}$$
 であり、したがって$\rho_\al=|\det T'| \rho_\be$すなわち$T_\sharp \al = \be$が回復される。ここで$T$は凸関数$\psi:x\mapsto \tfrac{1}{2}\dotp{x-\mean_\al}{A (x-\mean_\al)} + \dotp{\mean_\be}{x}$の勾配であることに注目し、\citeauthor{Brenier91}の定理~\citeyearpar{Brenier91}（注意~\ref{rem-exist-mongemap}を参照）を用いて$T$が最適であると結論する。写像$T$と対応するポテンシャル$\psi$の両方が図~\ref{fig-gaussians-2d-T}と図~\ref{fig-gaussians-2d-psi}に示されている。

$\rho_\al$の1次および2次モーメントに関する追加の計算により、この写像の輸送コストは次のようになる。
\eql{\label{eq-dist-gauss}
	\Wass_2^2( \al,\be ) = \norm{ \mean_\al - \mean_\be }^2 + \Bb(\cov_\al,\cov_\be)^2,
}
ここで$\Bb$はいわゆる\citeauthor{bures1969extension}計量~\citeyearpar{bures1969extension}であり、正定値行列間で定義される（\citet{,forrester2016relating}も参照）。
\eql{\label{eq-bure-defn}
	\Bb(\cov_\al,\cov_\be)^2 \eqdef \tr\pa{
		\cov_\al + \cov_\be - 2 ( \cov_\al^{1/2} \cov_\be \cov_\al^{1/2} )^{1/2}
	},
}
ここで$\cov^{1/2}$は行列の平方根である。$\Bb$は共分散行列上の距離であり、$\Bb^2$は両引数に関して凸であることを示せる。\todoK{ref ? proof ? }
$\cov_\al = \diag(r_i)_i$および$\cov_\be = \diag(s_i)_i$が対角行列の場合、ビューレス計量はヘリンジャー距離である。
\eq{
	\Bb(\cov_\al,\cov_\be) = \norm{ \sqrt{r}-\sqrt{s} }_2.
}
1次元ガウス分布の場合、$\Wass_2$はガウス分布の平均と標準偏差$(\mean,\sqrt{\cov})$をプロットした2次元平面上のユークリッド距離となる。図~\ref{fig-1d-gaussian}に示すとおりである。
%
ガウス分布のワッサーシュタイン幾何学の詳細な扱いについては\citet{takatsu2011wasserstein}を参照されたい。ビューレス計量に関する追加の考察については、最近の文献~\citep{malago2018wasserstein,bhatia2018bures}を参照されたい。単語の確率的埋め込みを計算するためのこの計量の最近の応用については~\citep{NIPS2018_8226}を、カルマンフィルタリングのロバストな拡張を計算するための使用については~\citep{NIPS2018_8067}を、再生核ヒルベルト空間における共分散関数への適用については~\citep{NIPS2017_7149}を参照されたい。
\end{rem2}
%%%%%%
\begin{figure}[h!]
\centering
\includegraphics[width=.7\linewidth]{gaussians-2d/gaussians1.pdf}
\caption{\label{fig-gaussians-2d-T} 2つのガウス分布$\rho_\al$と$\rho_\be$。密度の等高線プロットを用いて表され、平均と分散行列はそれぞれ$\mean_\al=(-2,0),\cov_\al=\frac{1}{2}\left(1 -\tfrac{1}{2};-\tfrac{1}{2}  1\right)$および$\mean_\be=(3,1), \cov_\be=\left(2, \tfrac{1}{2}; \tfrac{1}{2}, 1\right)$である。矢印は平面上に取ったランダムな点$x$を始点とし、それらの点の対応する写像$T(x)=\mean_\be + A(x-\mean_\al)$を終点とする。}
\end{figure}



%%%%%%%
\begin{rem2}{楕円等高分布間の距離}\label{rem-dist-elliptic}
\citeauthor{gelbrich1990formula}は注意~\ref{rem-dist-gaussians}で提供されたものよりも一般的な結果を提供する。ガウス分布間のビューレス計量はより一般的に\emph{楕円等高分布}に拡張される~\citeyearpar{gelbrich1990formula}。
%
簡潔に言えば、まず、与えられた平均と共分散行列を持つ2つの測度に対して、これらのそれぞれのパラメータを持つ2つのガウス分布間の距離が2つの測度間のワッサーシュタイン距離の下界であることを示せる~\citep[Theorem 2.1]{gelbrich1990formula}。さらに、閉形式~\eqref{eq-dist-gauss}は楕円等高密度の族に拡張される。2つの密度$\rho_\al$と$\rho_\be$がそのような族に属する場合、すなわち$\rho_\al$と$\rho_\be$が平均と正定値パラメータを用いて任意の点$x$に対して次のように書ける場合、

\eq{\begin{aligned}\rho_\al(x) = \frac{1}{\sqrt{\det(\A)}}h(\dotp{x-\mean_\al}{\A^{-1}(x-\mean_\al)})\\ \rho_\be(x)= \frac{1}{\sqrt{\det(\B)}} h(\dotp{x-\mean_\be}{\B^{-1}(x-\mean_\be)}),\end{aligned}}
同一の非負値関数$h$を用いて、積分$$\int_{\RR^\dim} h(\dotp{x}{x})dx=1$$を満たすとき、それらの最適輸送写像も線形写像\eqref{eq-transport-Bures}であり、ワッサーシュタイン距離も式$\eqref{eq-dist-gauss}$で与えられる。ただし、ビューレス計量のスケーリングが生成関数$h$にのみ依存してわずかに異なる。例えば、そのスケーリングはガウス分布の場合（$h(t)=e^{-t/2}$）は$1$であり、楕円体上の一様分布の場合（$h$が$[0,1]$の指示関数）は$1/(\dim+2)$である。この結果は、楕円分布の共分散行列がその正定値パラメータの定数倍であること~\citep[Theo. 4(ii)]{gomez2003survey}と、楕円分布間のワッサーシュタイン距離がそれらの共分散行列間のビューレス距離の関数であること~\citep[Cor. 2.5]{gelbrich1990formula}から従う。
\end{rem2}
%%%%%%
\begin{figure}[h!]
\centering
\includegraphics[width=.85\linewidth]{gaussians-2d/gaussians2.pdf}
\caption{\protect\label{fig-gaussians-2d-psi} 図~\ref{fig-gaussians-2d-T}で定義されたものと同じガウス分布$\rho_\al$と$\rho_\be$。今回は曲面として表現されている。上の曲面は$T=\nabla \psi$を満たすブレニエポテンシャル$\psi$であり、加法定数の違いを除いて定義される（ここでは+50）。視覚的な目的のため、両方のガウス密度は100倍されている。}
\end{figure}

\begin{figure}[h!]
\centering
\begin{tabular}{@{}c@{\hspace{1mm}}c@{}}
\includegraphics[width=.35\linewidth]{gaussian-1d/interp-density}
\includegraphics[width=.25\linewidth]{gaussian-1d/gaussian-param}
\end{tabular}
\caption{\label{fig-1d-gaussian}
2つの1次元ガウス分布間の変位補間の計算。
%
ガウス密度を$\Gg_{m,\si}(x) \eqdef \frac{1}{\sqrt{2\pi}s}e^{-\frac{(x-m)^2}{2s^2}}$と記すと、補間$\Gg_{(1-t)m_0+t m_1,(1-t)\si_0+t \si_1}$を示す。
}
\end{figure}



\todoK{Maybe cover the case of translation / scaling of the same input. Explain that $\Wass$ is invariant to isometries. }
